\chapter*{Preface}
\addcontentsline{toc}{chapter}{Preface}

\vspace{0.8cm}

This thesis investigates spatiotemporal deep-learning models for the magnetic
winding-flux field of solar active regions, derived from Solar Dynamics
Observatory (SDO) Helioseismic and Magnetic Imager (HMI) observations. The work
was carried out in the Department of Physics, Indian Institute of Technology
(Banaras Hindu University), Varanasi, as part of the Integrated Dual Degree
(B.Tech.\ + M.Tech.) programme in Engineering Physics.

\noindent\textbf{Subject.}\quad
The subject is the design, training and honest evaluation of a Convolutional
LSTM (ConvLSTM) encoder--forecaster that predicts the short-term evolution of
the two-dimensional winding-flux map of an active region, as a step toward
structure-aware solar flare forecasting.

\noindent\textbf{Scope.}\quad
The thesis is organised as three chapters. Chapter~\ref{ch:intro} introduces
solar flares, the space-weather forecasting problem, and the current
data-driven landscape, and identifies the gap addressed here: existing work
integrates the winding-flux map to a single scalar, discarding its spatial
structure. Chapter~\ref{ch:arch} specifies the data representation, the ConvLSTM
architecture, and the leave-one-active-region-out training procedure.
Chapter~\ref{ch:results} reports the forecasting results and analyses them
against a persistence baseline, distinguishing genuine predictive skill from the
persistence afforded by a smoothly evolving field.

\noindent\textbf{A note on honesty.}\quad
Because the accumulated total winding is strongly persistent, every forecast is
reported against a persistence baseline, and the effect of the spatial low-pass
is stated plainly in the text rather than obscured. No result in this thesis is
presented as stronger than the evidence supports.

\vspace{2cm}

\noindent
\hfill \theauthorname\\
\hfill Varanasi
